\subsection{Analytic soft-endpoint terms in the inner energy integral}
The $UV$ logarithmic sector illustrates why the regulator must be
kept until both endpoints have been integrated. Put
\begin{equation}
 k=\frac{(v-1)^2}{Hv},\qquad u=\frac{x}{H},\qquad
 y=\frac{H(1-u)}{1-ku},\qquad
 (xy)^\eta=H^{2\eta}u^\eta(1-u)^\eta(1-ku)^{-\eta}.
\end{equation}
For $1<v<s=1+H$, $0<k<H/s<1$. Thus the third pole
$x=H/k$ is outside the entire inner interval $0<x<H$.
The rational coefficient of $\ln v$ has precisely the form
\begin{equation}
 R(x,v)=C_0+\frac{C_{01}}x+\frac{C_{02}}{x^2}
 +\frac{C_{H1}}{x-H}+\frac{C_{H2}}{(x-H)^2}
 +\frac{C_{a1}}{x-H/k}+\frac{C_{a2}}{(x-H/k)^2}.
\end{equation}
The coefficients are obtained by ordinary partial fractions of the
explicit $R_{UV}$ above after the change of variables.
The script \nolinkurl{compton_real_inner_analytic.py} performs this
algebra and verifies exact reconstruction of the original rational
function. No kinematic fitting is involved.

Each integral is an Euler beta function times a Gauss
hypergeometric function. For example, for sufficiently large
$\operatorname{Re}\eta$,
\begin{align}
 \int_0^H\frac{(xy)^\eta}{x^n}\dd x
 &=H^{1-n+2\eta}B(1+\eta-n,1+\eta)
 {}_2F_1(\eta,1+\eta-n;2+2\eta-n;k),\\
 \int_0^H\frac{(xy)^\eta}{(x-H)^n}\dd x
 &=(-1)^nH^{1-n+2\eta}B(1+\eta,1+\eta-n)
 {}_2F_1(\eta,1+\eta;2+2\eta-n;k),\\
 \int_0^H\frac{(xy)^\eta}{(x-H/k)^n}\dd x
 &=(-H/k)^{-n}H^{1+2\eta}B(1+\eta,1+\eta)
 {}_2F_1(n+\eta,1+\eta;2+2\eta;k).
\end{align}
Here $B(a,b)=\Gamma(a)\Gamma(b)/\Gamma(a+b)$. The Gauss function is
\begin{equation}
 {}_2F_1(a,b;c;k)=\sum_{n=0}^{\infty}
 \frac{(a)_n(b)_n}{(c)_n n!}k^n,\qquad |k|<1,
 \qquad (a)_n=\frac{\Gamma(a+n)}{\Gamma(a)},\quad (a)_0=1.
\end{equation}
Expanding $(1-ku)^{-a}$ and integrating each power against
$u^{b-1}(1-u)^{c-b-1}$ gives the displayed beta integrals in their
convergence domain. These expressions
then uniquely define the continuation to $\eta=0$. Their Laurent
expansion gives, at fixed $v$ in the interior,
\begin{align}
 \int_0^H (xy)^\eta R(x,v)\dd x
 &=\frac{C_{01}-C_{H1}}\eta+F_{UV}(v)+\Order(\eta),\\
 F_{UV}(v)&=H C_0+\frac{k-2}{H}C_{02}
 -\frac{2-k}{H(1-k)}C_{H2}
 +2(C_{01}-C_{H1})\ln H\\
 &\quad +(C_{H1}+C_{a1})\ln(1-k)
 +\frac{k^2}{H(1-k)}C_{a2}.
 \label{eq:comp-real-UV-inner-finite}
\end{align}
The double-pole terms have finite limits but no $1/\eta$ pole.
Their finite constants cannot be obtained by simply setting
$\eta=0$ in the integrand. For example,
\begin{equation}
 (1-u)^\eta(1-ku)^{-\eta}
 =1+\eta(k-1)u+\Order(u^2),\qquad
 \int_0^1u^{\eta-2}\dd u=\frac1{\eta-1}.
\end{equation}
The linear term contributes $\eta(k-1)/\eta=k-1$, which combines
with $-1$ from the first term to give $k-2$.
This explicitly accounts for the finite endpoint contribution.

Equation~\eqref{eq:comp-real-UV-inner-finite} still has to be
combined with the local $z=0$ subtraction and integrated over $v$.
It must not be integrated as an isolated term across a spurious
Gram pole before those cancellations have been made.
